\documentclass[11pt]{article}
\usepackage[a4paper,margin=23mm]{geometry}
\usepackage[T1]{fontenc}
\usepackage{lmodern,amsmath,amssymb,amsthm,mathtools,microtype,booktabs}
\usepackage[colorlinks=true,linkcolor=blue,citecolor=blue,urlcolor=blue]{hyperref}
\newcommand{\E}{\mathbb E}
\newcommand{\Tr}{\operatorname{Tr}}
\newcommand{\tr}{\operatorname{tr}}
\newcommand{\HH}{\mathcal H}
\newcommand{\norm}[1]{\left\lVert#1\right\rVert}
\newcommand{\dd}{\mathrm d}
\newtheorem{theorem}{Theorem}
\title{Same-disorder curvature in the Slater selected algebra\\
\large Fixed-word source calculus, central pinching and covariant observables}
\author{Alan Day \and Chengheng\thanks{AI research assistant (OpenAI).}}
\date{8 November 2026\quad Local research draft}
\begin{document}\maketitle
\begin{abstract}
We construct all fixed mixed moments of the actual quadratic curvature
field with the Hamiltonian, local size fluctuation and Slater-selected
operators. The construction retains the original Gaussian disorder.
Separate empty and occupied local matrix elements extend the existing
nilpotent loss-source calculus to selected chords. Without selected
chords the source exponent agrees exactly with the prior tracial one,
which proves the central-sector moment and pinching statements directly.
The polynomial orbit-density estimate transfers these errors to arbitrary
fixed selected polynomial preparations. It also makes the covariant
three-field algebra a bounded normal representation on the enlarged
GNS space. A nonzero curvature/double-selected overlap rules out an
independent-probe substitution in this state.
\end{abstract}

\section{Microscopic definitions and the result}
Keep standard CAR, overlapping Hermitian complex Gaussian $p$-body
couplings, the same $J$ in every field, a fixed half-filled Slater vector
$O$, and $\varphi_N=\E\langle O|\,\cdot\,|O\rangle$. Let
$m=N/2$, $\lambda_N=p^2/N\to\lambda>0$, $q_N=e^{-4\lambda_N}$,
$C_Q=\E H^2$, and $C_m$ its scalar value on $\bigwedge^m\mathbb C^N$.
All limits fix the word and $\lambda$ first. Set
\begin{align}
 D X&=QX-\sum_i c_i^\dagger Xc_i,&
 \mathcal L X&=\frac{NX-\sum_i(c_i^\dagger Xc_i+c_iXc_i^\dagger)}2,\label{losses}\\
 Z_2&=H^2-C_Q,&
 \kappa_N&=2\lambda_N\frac{1-q_N}{1+q_N},\nonumber\\
 j_N&=C_m^{-1}(2p-D-\kappa_N)Z_2\big|_m,&
 v_N&=4\lambda_N+\frac{16\lambda_N^2q_N}{1+q_N},\label{jdef}\\
 h_N&=H_m/\sqrt{C_m},&
 W_N&=(p-\mathcal L-\lambda_N)H\big|_m/\sqrt{C_m},\nonumber\\
 u_N&=W_N/\sqrt{\lambda_N},&x_N&=j_N/\sqrt{v_N}.\nonumber
\end{align}
The loss maps act on full-Fock operators \emph{before} sector restriction.
The deterministic $C_m^{-1}$ is applied afterwards; replacing $H$ by
the sector-dependent $H/\sqrt{C_Q}$ before applying $D$ would change
the field. $M_N^\dagger$ is the original normalized occupied-to-empty
$p$-particle creation block, with $\varphi_N(M_NM_N^\dagger)=1$.
These are dimensionless string-energy observables, with $\hbar=1$.

\begin{theorem}\label{main}
Every fixed word in $h_N,W_N,j_N,M_N,M_N^\dagger$ has the limit
specified by the finite diagram rule in Sections 2--3. The limiting
functional is positive. Its covariant restriction $(h,u,x)$ has the
existing tracial mixed-$q$ law with lengths $(1,1,2)$. Exact energy
pinching of $x_N$ vanishes in every fixed central Schatten norm and
on every fixed polynomial preparation in this full family.
\end{theorem}

\newpage
\section{A local source rule with the selected state retained}
Expand $j$ into two adjacent Hamiltonian slots with a whole-block $D$
source; expand $W$ into one slot with an $\mathcal L$ source. Gaussian
Wick centering of $Z_2$ forbids pairs internal to its two-slot block.
Other Hamiltonian slots are ordinary. Write a Gaussian pairing as
$\pi=\{(a_e,b_e):a_e<b_e\}_{e=1}^k$. At the right endpoint use
$C_{P_e}^\dagger C_{Q_e}$, and at the left its adjoint. An unrestricted
edge is $G$. A forward selected edge is $F$, requiring $P_e$ empty and
$Q_e$ occupied in $O$. A backward selected edge has the reverse
constraint. Two creations or two annihilations cannot pair.

Use one nilpotent variable $z_a$, $z_a^2=0$, for each marked block.
For a one-site matrix $X=\begin{psmallmatrix}a&b\\c&d\end{psmallmatrix}$,
the local maps are
\begin{equation}\label{localmaps}
 D_iX=\begin{pmatrix}0&b/2\\c/2&d-a\end{pmatrix},\qquad
 \mathcal L_iX=\frac12\begin{pmatrix}a-d&b\\c&d-a\end{pmatrix}.
\end{equation}
For membership $S\subset\{P_e,Q_e\}_{e=1}^k$, multiply the local
CAR matrices inside each block, apply $1-z_aD_i$ or
$1-z_a\mathcal L_i$, and multiply the blocks in their original order.
Take the $(o,o)$ entry for $o=0,1$, rather than the local trace. Multiply
by the Jordan--Wigner regrouping sign
\begin{equation}\label{sign}
 (-1)^{\sum_{e<f:\,e\text{ crosses }f} t_e(S)t_f(S)},\qquad
 t_e(S)=|S\cap\{P_e,Q_e\}|.
\end{equation}
Set the answer to zero for any violated selected-label restriction.
Call the resulting polynomial $f_o(S;z)$.

Each label has one correct singleton occupancy: $P_e$ at $o=0$,
$Q_e$ at $o=1$. For a forward diagram its singleton polynomial
$s_i(z)$ has constant term one. Because no marked block contains
both endpoints of a Gaussian pair, each incident singleton contributes
$1-z_a/2$ to that block. For a block of $d_a$ slots,
\begin{equation}\label{cancel}
 e_{\rm nil}^{p\sum_a d_a z_a}\prod_i s_i(z)^p=1.
\end{equation}
This removes the extensive terms of $pd_a-D$ or
$pd_a-\mathcal L$ in all source coefficients, not just first
derivatives. Define normalized collision polynomials
\begin{equation}\label{ratio}
 F_o(S;z)=f_o(S;z)\prod_{i\in S}s_i(z)^{-1}.
\end{equation}
The inverse exists in the finite nilpotent algebra. In particular,
we do not divide by a possibly vanishing collision matrix element.
The local matrix products and the sign are the finite CAR rule;
no Poisson approximation has yet been made.

\newpage
\section{The fixed-word limit and its domination}
Let $b(\pi)$ be the number of Hamiltonian--selected edges and let
$\nu(\pi)$ count crossings plus nested pairs whose inner edge is
selected. Only compatible forward diagrams survive. Define
\begin{equation}\label{K}
 K_\pi(z)=2\sum_{i<j}\sum_{o=0,1}F_o(\{i,j\};z)
                 -2k(k-1)-2k_G,
 \quad K_\pi(0)=-4\nu(\pi),
\end{equation}
where $k_G$ is the number of unrestricted edges. Put
$q=e^{-4\lambda}$, $\gamma=e^{-\lambda}=q^{1/4}$ and
$\kappa=2\lambda(1-q)/(1+q)$. For $r$ marks, the diagram contribution is
\begin{equation}\label{rule}
 \gamma^{b(\pi)}[z_1\cdots z_r]
 \exp_{\rm nil}\left\{\lambda K_\pi(z)
       -\lambda\sum_{a:W}z_a-\kappa\sum_{a:j}z_a\right\}.
\end{equation}
The constant exponential is $q^{\nu(\pi)}$; nilpotence applies to
the remaining terms. Sum \eqref{rule} over the allowed pairings.
For normalized $u,x$, divide additionally by
$\lambda^{n_W/2}v^{n_j/2}$, with $v=\lim v_N$.

Here is a fixed-degree convergence argument, including signed terms.
Let $L_N=\binom mp^2$ and expand the exact coefficient of the empty
and occupied site polynomials $f_0^mf_1^m$. After \eqref{cancel},
assign $n_{o,S}$ collision sites of membership $|S|\ge2$, with
$d_i=\sum_{o,S\ni i}n_{o,S}$. The normalized combinatorial factor is
\begin{equation}\label{defect}
 \frac{(m)_{U_0}(m)_{U_1}}{(m)_p^{2k}}
 \prod_i(p)_{d_i}\prod_{o,S}
                   \frac{F_o(S;z)^{n_{o,S}}}{n_{o,S}!},
\end{equation}
where $U_o$ counts the singleton and collision sites in occupancy $o$.
The source-free counting is unchanged from the previous Slater proof.
The coefficient norm of each $F_o$ is bounded by a fixed $B$.
Its falling-factorial estimate therefore bounds the absolute series
by a constant times the product of exponential series with activities
\begin{equation}
 B\alpha_{S,N},\qquad \alpha_{S,N}=m(p/m)^{|S|}.
\end{equation}
For fixed slots and marks their total is bounded. The $|S|\ge3$
activity tends to zero, and the total-defect tail above $R$ tends to
zero uniformly as $R\to\infty$. This establishes domination before
discarding higher intersections or taking source coefficients.

At fixed defects, the falling-factorial background is
$e^{-2\lambda k(k-1)}$ and each double membership has activity
$p^2/m\to2\lambda$. The Hamiltonian normalization uses
$D_N=\binom mp\binom{m+p}p$, with $D_N/L_N\to e^{2\lambda}$.
There are $n_H=2k_G+b$ Hamiltonian slots, giving the additional
$e^{-\lambda n_H}$. These factors produce \eqref{K}--\eqref{rule}.
For a backward selected edge a correct singleton is forbidden,
forcing $d_i=p$ and hence a divergent number of defects. Its
contribution vanishes by the same tail bound. The source maps
preserve singleton parity and cannot restore that forbidden entry.
All fixed moments, including those of $P_N^\dagger P_N$, converge
and stay bounded. Their finite-$N$ positivity passes to the limit.

\newpage
\section{The central-sector identification is a direct calculation}
When all edges are $G$, the empty and occupied source polynomials
satisfy $T(S;z)=(f_0(S;z)+f_1(S;z))/2$, where $T$ is the old local
normalized trace. Its singleton is $t_i=s_i/2$. Thus its normalized
two-label collision ratio is
\begin{equation}\label{central}
 R_{ij}^{\rm trace}(z)=\frac{T(\{i,j\};z)}{t_i(z)t_j(z)}
   =2\sum_{o=0,1}F_o(\{i,j\};z).
\end{equation}
The tracial background is $-\binom{2k}{2}-k=-2k^2$.
For all $G$, the background in \eqref{K} is also $-2k^2$.
Therefore the entire source exponent, with every coefficient
retained, is identical to the previous tracial one. This proves
the central limit from its own finite coefficient expansion; it is
not an inference from concentration of the full-Fock charge average.

For a covariant word, $\E$ of that word is scalar on
$\bigwedge^m\mathbb C^N$. Its Slater expectation consequently equals
its normalized sector trace exactly. The established endpoint
cancellation of the whole two-slot $D$ field and the one-slot
$\mathcal L$ rule now apply to \eqref{central}. They give the mixed-$q$
tracial law for $(h,u,x)$ with covariance the identity and crossing matrix
\begin{equation}\label{mixedq}
 \mathsf Q=\begin{pmatrix}
 q&q&q^2\\q&q&q^2\\q^2&q^2&q^4
 \end{pmatrix}.
\end{equation}
This tracial law and its usual bounded representation were already
used in the full-trace bilocal stage. What is new here is its direct
central-sector derivation and its coupling to selected observables.
In particular all central fixed Schatten moments of $x_N$ are bounded.

Let $\mu_q$ be the unit-variance $q$-Hermite measure and $H_n$ its
monic polynomials, with squared norm $[n]_q!$. The two-energy measure
of $x_N$ converges to the established Mehler measure
\begin{equation}\label{measure}
 \dd\nu(E,E')=K_{q^2}(E,E')\,\dd\mu_q(E)\,\dd\mu_q(E'),\qquad
 K_r(E,E')=\sum_{n\ge0}\frac{r^nH_n(E)H_n(E')}{[n]_q!}.
\end{equation}
The finite measure is positive:
$\nu_N=\E\tr_m\sum_{a,b}\Pi_a x_N\Pi_bx_N\,
\delta_{(E_a,E_b)}$. All its polynomial moments follow from the
central mixed-word limit. Higher energy moments provide tightness
and uniform integrability; the compact limiting support determines
the measure. Since $q^2<1$, \eqref{measure} is absolutely continuous
with respect to the product of the nonatomic $q$-Hermite measures.
It has no mass on $E=E'$.

Portmanteau on the closed diagonal therefore gives
\begin{equation}\label{pinch}
 \norm{\mathcal P_{H_N}x_N}_{2,m}^2=\nu_N(E=E')\longrightarrow0.
\end{equation}
Pinching is contractive in every Schatten norm. Interpolating
\eqref{pinch} with higher fixed central moments gives convergence
to zero in every fixed finite $L^r(\tr_m)$, including all degenerate
energy blocks. No long-time and large-$N$ limits are interchanged.

\newpage
\section{Actual curvature and selected polynomial preparations}
The previously established finite-$N$ identity is
\begin{equation}\label{actual}
 A_{0,N}=(1-\mathcal P_{H_N})\bigl((2p-D)Z_2\bigr)\big|_m.
\end{equation}
The omitted regression term in \eqref{jdef} is a function of $H_m$
plus a scalar and is killed by $1-\mathcal P_{H_N}$. Thus, exactly,
\begin{equation}\label{an}
 a_N:=\frac{A_{0,N}}{C_m\sqrt{v_N}}
       =(1-\mathcal P_{H_N})x_N.
\end{equation}
This $a_N$ denotes a normalized observable, not the preparation norm
used in the preceding orbit-density draft. The physical $J_{2,N}$
defined there with the limiting regression coefficient differs from
$C_mj_N$ by a coefficient tending to zero times $Z_2$, and a scalar
proportional to $C-C_m$. Both differences vanish in every fixed
central Schatten norm. Equation \eqref{an} itself needs no such
asymptotic replacement.

Now allow $P_N$ to be any fixed polynomial in
$h_N,u_N,x_N,M_N,M_N^\dagger$. It has a fixed Gaussian degree,
is stabilizer-covariant, and has bounded preparation norm by the
source theorem. The preceding polynomial orbit-density result yields
a constant $C_P$, independent of $N$, with
\begin{equation}\label{density}
 \norm{R_{P,N}}_{L^2(\tr_m)}\le C_P,\qquad
 \E\norm{X_NP_NO}^2\le C_P\norm{X_N}_{4,m}^2
\end{equation}
for every covariant $X_N$. Combining \eqref{pinch} and \eqref{an},
\begin{equation}\label{selectedpinch}
 \E\norm{(a_N-x_N)P_NO}^2\longrightarrow0.
\end{equation}
The actual normalized neutral slope $b_N$ also obeys
$\E\norm{(b_N-u_N)P_NO}^2\to0$ for this enlarged class of preparations:
its already proved central tracial error is unchanged, and the
new $P_N$ now has the required bounded orbit density.

More generally, let $W_N$ be a fixed covariant block built from
$h_N,a_N,b_N$ and fixed $h_N$-time phases. Tracial H\"older and
telescoping first replace $a_N,b_N$ by $x_N,u_N$ in $L^2$.
Equation \eqref{density} and polarization of $P_N,Q_N$ then give
\begin{equation}\label{sandwich}
 \varphi_N\bigl(P_N^\dagger W_N(h_N,a_N,b_N)Q_N\bigr)
 \longrightarrow\varphi\bigl(P^\dagger W(h,x,u)Q\bigr).
\end{equation}
For the time phases, central Gaussian moments of $h_N$ give a
uniform Schatten Taylor remainder of order
$S^n[(rn-1)!!]^{1/r}/n!$ at fixed even $r$.
Each fixed number of covariant insertions has uniformly bounded higher
norms, so H\"older gives $L^2$ truncation uniformly on compact real
time windows. Apply the density bound before removing the truncation.
This proves the time version of \eqref{sandwich}.

The selected factors in $P,Q$ are unrestricted fixed polynomials.
Several nonpolynomial exact-pinched blocks separated by selected
factors still require an additional argument: their intervening
preparations are not Gaussian polynomials. The present result does
not replace that problem by convergence on a core.

\newpage
\section{Bounded covariant observables on the joint GNS space}
The source theorem defines a positive functional on the full
polynomial $*$-algebra. Form its GNS quotient and completion;
let $\mathcal C$ be its dense polynomial core. The fields
$h,u,x$ are symmetric there, while $M^\dagger$ is the algebraic
adjoint of $M$. We can identify the covariant action without
guessing an independent selected representation.

Let $\mathcal D_0$ be the tracial mixed-$q$ polynomial algebra in
\eqref{mixedq}, with faithful trace $\tau$, and let $A=A^*$ be one
of its polynomials. For $\psi=P\Omega$, the finite-$N$ density bound
and both moment limits imply, for every fixed integer $n$,
\begin{equation}\label{power}
 \norm{A^n\psi}^2
 \le C_P\left[\tau(A^{4n})\right]^{1/2}
 \le C_P\norm{A}_{\mathcal D}^{2n}.
\end{equation}
Here $\mathcal D$ is the usual tracial $C^*$ completion. The constant
$C_P$ does not grow with $n$; it is supplied by the fixed degree of
the preparation, not the degree of the observable.

Log-convexity of the positive polynomial moments, or repeated
Cauchy--Schwarz for $n=2^k$, gives
\begin{equation}
 \norm{A\psi}^2\le
 \norm{\psi}^{2-2/n}\norm{A^n\psi}^{2/n}
 \le\norm{\psi}^{2-2/n}C_P^{1/n}\norm{A}_{\mathcal D}^2.
\end{equation}
Letting $n\to\infty$ proves
$\norm{A\psi}\le\norm{A}_{\mathcal D}\norm{\psi}$.
Thus covariant null polynomials act as zero, all these operators extend
boundedly, and the representation extends to $\mathcal D$.
The vacuum restriction is $\tau$, so it is faithful. In particular,
\begin{equation}
 \norm{h},\norm{u}\le\frac2{\sqrt{1-q}},\qquad
 \norm{x}\le\frac2{\sqrt{1-q^4}}.
\end{equation}

There is also a useful normality statement. For every core vector,
\begin{equation}\label{normal}
 |\langle\psi,A\psi\rangle|\le C_P\norm{A}_{L^2(\tau)}.
\end{equation}
Riesz representation gives an $L^2(\tau)$ density $\rho_\psi$ for
this functional. Positivity on polynomial squares extends by bounded
approximation to positivity of the density; its trace is
$\norm{\psi}^2$. Hence these vector functionals are normal on the
tracial von Neumann completion of $\mathcal D$. Polarization and
density of $\mathcal C$ give a normal representation there.

This establishes a common representation of the actual covariant
fields and the selected polynomial algebra. It does not classify
the complete enlarged algebra, prove compactness of $M$ on this
larger space, or provide a closed open-chord transfer for $j$.
Compactness in the older $h,u,M$ cyclic subspace alone would not
prove any of those statements.

\newpage
\section{Same-disorder mixed moments that distinguish the construction}
The diagram rule gives the following unnormalized $j$ moments:
\begin{align}
 \varphi(j^2)&=v,&\varphi(jh^2)&=0,\label{tests1}\\
 \varphi\bigl(j(M^\dagger)^2\bigr)&=-2\kappa\gamma^2q,&
 \varphi(jM^\dagger h)&=-2\kappa\gamma q,\label{tests2}\\
 \varphi(jhM^\dagger)&=0,&
 \varphi(jWM^\dagger)&=\varphi(jM^\dagger W)=0,\nonumber\\
 \varphi(MjM^\dagger)&=0,&
 \varphi(Mj^2M^\dagger)&=v.\label{tests3}
\end{align}
The order in \eqref{tests2} matters. To see its first identity directly,
pair the two adjacent $j$ slots with the two selected creations.
Every surviving label assignment has disjoint occupied-to-empty
hops. The unregressed loss deficit $2p-D$ on that pure $2p$-body
hop is zero. The remaining regression term is
$-\kappa\varphi(h^2(M^\dagger)^2)$, and the two forward pairings give
$\varphi(h^2(M^\dagger)^2)=2\gamma^2q$.
This finite-label reasoning agrees with the normalized source rule.

Since $\kappa>0$ for every fixed $\lambda>0$, the overlap is
nonzero on the physical curve $\gamma=e^{-\lambda}$. A separately
centered, disorder-independent probe would have zero such overlap.
The tracial marginal law in \eqref{mixedq} therefore cannot be used
as a substitute for the common-disorder selected moments.

Equation \eqref{sandwich} transfers these identities to actual
curvature insertions. For example,
\begin{equation}\label{actualtests}
 \lim_N\varphi_N\bigl(a_N(M_N^\dagger)^2\bigr)
   =-\frac{2\kappa\gamma^2q}{\sqrt v},\qquad
 \lim_N\varphi_N(M_Na_N^2M_N^\dagger)=1.
\end{equation}
The second preparation already has unit norm. The first expression
is an off-diagonal matrix element, not an expectation in the
normalized double-excited state. Its adjoint gives the reversed
matrix element. The physical curvature is recovered by multiplying
by $C_m\sqrt{v_N}$ before taking the limit.

For an algebraic calibration, form the Gram matrix of the fourteen
preparations
\begin{equation}
 1,\ H,\ W,\ M^\dagger,\ H^2,\ HW,\ WH,\ W^2,\
 M^\dagger H,\ HM^\dagger,\ M^\dagger W,\ WM^\dagger,\
 (M^\dagger)^2,\ j.
\end{equation}
Its exact LDL pivots are positive at the two saved rational choices
$(\gamma,\lambda)=(1/2,1/4),(2/3,1/9)$.
These are formal polynomial diagnostics, not points on
$\gamma=e^{-\lambda}$ and not a proof of the full infinite-dimensional
positivity or of a new orthogonal sector. General physical positivity
comes from the finite-$N$ construction and its moment limit.

\newpage
\section{Checks, sources and the next interface}
The independent finite checker builds full-Fock integer CAR matrices
for $(N,p)=(2,1),(4,1),(4,2)$. It applies the global maps in
\eqref{losses} to entire blocks and compares every source coefficient
with the product of the empty/occupied local polynomials. It includes
forward, backward and unrestricted label restrictions, internal
Gaussian pairs at this finite algebra level, and mixed $D/\mathcal L$
block layouts. The limiting driver separately enforces Wick exclusions.
There are 592,704 exact comparisons, including 40 negative contributions.
The maximum recorded integer magnitude is 16; the bounded-size matrix
products are safely within integer range.

Further checks compare 326 all-$G$ source exponents with the existing
full-trace helper, their words with the established mixed-field
formula, and 340 words without $j$ with the older Slater $H/U/M$
transfer. Ten symbolic mixed expressions include
\eqref{tests1}--\eqref{tests3}; two fourteen-vector Gram checks use
exact rational arithmetic. These finite comparisons test signs,
source placement, counterterms and normalization. The arbitrary
fixed-word convergence, bounded representation and exact pinching
are proved above, not inferred from their finite coverage.

The source method builds on our already recorded full-trace
multimark and bilocal derivations, and on the fixed-Slater defect
expansion. The physical fixed-charge problem and the complex-SYK
CAR/chord conventions have established precedents \cite{bnr}.
The dependence of the emergent algebra on retained state-adapted
observables is also explicit in recent pure-state SYK work
\cite{rx}; its Majorana/KM construction is not substituted for the
present charge-conserving same-$J$ composite calculation.

The next task is to turn the constructive selected-composite rule
into an efficient transfer and identify its low-lying modules,
then include the higher charged and neutral descendants of the
full bilocal. The completed full-trace finite-time bilocal four-point
construction remains available. The generalized source rule can
treat fixed higher whole-block Wick marks, but their actual rank
and charge-removal interfaces, uniform time tails in the selected
state, and interlaced exact-pinching moments require separate work.
The bulk spatial and mass generators and finite-$G$ dynamics remain
open. This stage is a microscopic observable construction, not a
completed de Sitter duality.

\begin{thebibliography}{9}
\bibitem{bnr} M. Berkooz, V. Narovlansky and H. Raj,
\emph{Complex Sachdev--Ye--Kitaev model in the double scaling limit},
arXiv:2006.13983v2, especially Sections 3 and 5.
\url{https://arxiv.org/html/2006.13983v2}.
\bibitem{rx} H. Rajgadia and J. Xu,
\emph{Emergent States and Algebras from the Double-Scaling limit
of Pure States in SYK}, arXiv:2604.14387v3, Sections 3--4.
\url{https://arxiv.org/html/2604.14387v3}.
\bibitem{local} A. Day and Chengheng, local project stages:
\emph{Slater joint words}, \emph{Slater covariant},
\emph{Slater polynomial density}, \emph{Singlet multimark}, and
\emph{Bilocal four point}. Their manuscripts, pure source helpers,
scientific outputs and provenance are retained in this repository.
\end{thebibliography}
\end{document}
